\documentclass[letterpaper,12pt]{article}
\usepackage[paperheight=27.94cm,paperwidth=21.59cm,bindingoffset=0in,left=3cm,right=2.0cm,top=3.5cm,bottom=2.5cm,headheight=15pt,headsep=1.0\baselineskip]{geometry}
\usepackage{graphicx}
\usepackage{amsmath,amssymb}
\usepackage[spanish,es-tabla]{babel}
\usepackage{booktabs}
\usepackage{tabularx}
\usepackage{float}
\usepackage{cite}
\usepackage[utf8]{inputenc}
\usepackage{fancyhdr}
\usepackage{xurl}
\usepackage{hyperref}
\usepackage{listings}
\usepackage{xcolor}

% Configuración del formato para el código VHDL
\lstset{
    language=VHDL,
    basicstyle=\ttfamily\footnotesize,
    keywordstyle=\color{blue}\bfseries,
    commentstyle=\color{green!60!black},
    stringstyle=\color{red},
    numbers=left,
    numberstyle=\tiny\color{gray},
    stepnumber=1,
    numbersep=10pt,
    backgroundcolor=\color{black!5},
    frame=single,
    rulecolor=\color{black!30},
    breaklines=true,
    breakatwhitespace=false,
    showstringspaces=false,
    captionpos=b,
    inputencoding=utf8,
    extendedchars=true,
    literate=
        {á}{{\'a}}1 {é}{{\'e}}1 {í}{{\'i}}1 {ó}{{\'o}}1 {ú}{{\'u}}1
        {ñ}{{\~n}}1 {Á}{{\'A}}1 {É}{{\'E}}1 {Í}{{\'I}}1 {Ó}{{\'O}}1
        {Ú}{{\'U}}1 {Ñ}{{\~N}}1 {¡}{{\textexclamdown}}1 {¿}{{\textquestiondown}}1
}

\renewcommand{\tablename}{Tabla}
\pagestyle{fancy}
\fancyhf{}
\rhead{\thepage}
\lhead{Laboratorio 3: Máquina de estados finitos para control de acceso digital}

\begin{document}

% ==============================
% 1. CARÁTULA
% ==============================
\begin{titlepage}
    \centering
    \includegraphics[width=0.55\textwidth]{Logo_Universidad_Diego_Portales.png}\\[1.2 cm]
    
    \textsc{\large Facultad de Ingeniería y Ciencias}\\[0.2 cm]
    \textsc{\large Escuela de Informática y Telecomunicaciones}\\[1.5 cm]
    
    \textsc{\Large Arquitectura y Organización de Computadores}\\[0.5 cm]
    
    \rule{\linewidth}{0.5 mm} \\[0.4 cm]
    {\huge \bfseries Laboratorio 3:\\[0.3cm] \Large Máquina de estados finitos para control de acceso digital}\\
    \rule{\linewidth}{0.5 mm} \\[2.0 cm]
    
    \large \emph{Integrantes:}\\[0.3 cm]
    {\Large Benjamín Oyarzún}\\[0.2cm]
    {\Large Lucas Herrada}\\[0.2cm]
    {\Large Joakin Mac-Auliffe}

    \vfill
    {\large Noviembre de 2026}
\end{titlepage}

\newpage
\tableofcontents
\newpage

% ==============================
% 2. INTRODUCCIÓN
% ==============================
\section{Introducción}
El presente informe documenta el proceso de diseño, simulación e implementación en hardware de un sistema de control de acceso digital. Para lograr este cometido, se ha utilizado una arquitectura basada en una Máquina de Estados Finitos (FSM) modelada mediante el lenguaje de descripción de hardware VHDL. 

Los sistemas secuenciales y las máquinas de estados constituyen el núcleo de las unidades de control en la arquitectura de computadores. Este proyecto permite poner en práctica dichos conceptos mediante el despliegue del diseño sobre una placa FPGA Nexys 4 DDR Artix-7, abordando desafíos inherentes al hardware real, como la sincronización de señales externas y la generación de pulsos de un ciclo de reloj a partir de pulsadores.

% ==============================
% 3. OBJETIVOS
% ==============================
\section{Objetivos}
El objetivo principal de este laboratorio es diseñar, simular e implementar en una FPGA una máquina de estados finitos (FSM) para controlar un sistema de acceso mediante una clave binaria ingresada secuencialmente.

Para alcanzar este propósito, el sistema debe cumplir con los siguientes objetivos específicos:
\begin{itemize}
    \item Detectar de forma exacta una clave correcta predefinida (10110).
    \item Indicar visualmente si se ha ingresado una clave incorrecta.
    \item Contabilizar el número de intentos fallidos mediante un registro interno.
    \item Bloquear permanentemente el sistema después de registrar tres errores.
    \item Permitir el reinicio total del sistema y contadores exclusivamente mediante una señal de reset global.
\end{itemize}

% ==============================
% 4. FUNDAMENTOS TEÓRICOS
% ==============================
\section{Fundamentos teóricos}
El desarrollo de este sistema se sustenta en los siguientes conceptos de lógica digital:
\begin{itemize}
    \item \textbf{Circuitos secuenciales:} A diferencia de la lógica combinacional, incorporan elementos de memoria. Sus salidas dependen tanto de las entradas actuales como de los estados previos.
    \item \textbf{Máquina de Moore:} Es un tipo de máquina de estados donde las salidas dependen única y exclusivamente del estado actual en el que se encuentra el sistema, simplificando la lógica de control visual.
    \item \textbf{Reloj (CLK):} Consiste en una señal oscilante que rige el tiempo del sistema. En este diseño, los cambios de estado ocurren únicamente durante el flanco ascendente del reloj.
\end{itemize}

% ==============================
% 5. DESCRIPCIÓN FUNCIONAL
% ==============================
\section{Descripción funcional}
El sistema controla el acceso validando la secuencia binaria \texttt{10110}. El usuario introduce cada bit utilizando el interruptor físico \texttt{SW0} y confirma la entrada presionando el botón \texttt{BTNC} (señal ENTER).

Para proteger la seguridad de la clave, la máquina de estados se ha diseñado con un retardo en la notificación de error: si el usuario se equivoca en cualquier punto, el sistema no lo alerta de inmediato, sino que se desvía a una rama de estados "sombra" que consume los bits restantes hasta alcanzar los cinco ingresos. Al finalizar la secuencia, si los bits fueron correctos se activa la señal \texttt{ACCESO\_OK} en el \texttt{LED1}. Si hubo errores, se enciende la señal \texttt{ERROR} en el \texttt{LED2} y se incrementa el contador de intentos fallidos.

Tras acumular tres errores, el sistema entra en un estado irreversible donde se enciende la señal \texttt{BLOQUEADO} (\texttt{LED3}), ignorando cualquier ingreso posterior. La única forma de salir de este bloqueo es presionando el botón \texttt{BTNU} (señal RESET).

% ==============================
% 6. DIAGRAMA DE BLOQUES
% ==============================
\section{Diagrama de bloques}
El diagrama de bloques representa la interconexión entre las entradas físicas (reloj de 100 MHz en \texttt{E3}, pulsador central \texttt{BTNC} como \texttt{ENTER}, pulsador superior \texttt{BTNU} como \texttt{RESET} y el interruptor \texttt{SW0} como dato \texttt{X}), el bloque acondicionador \texttt{detector\_flanco} que sincroniza y genera un pulso de un ciclo de reloj (\texttt{pulso\_enter}), la máquina de estados finitos tipo Moore \texttt{fsm\_acceso} con su contador de intentos integrado, y el mapeo final de los indicadores hacia los LEDs físicos (LED0 a LED9).

\begin{figure}[H]
    \centering
    \includegraphics[width=1\linewidth]{blockdiagram.png}
    \caption{Diagrama de bloques general del sistema implementado.}
    \label{fig:diagrama_bloques}
\end{figure}

% ==============================
% 7. DIAGRAMA DE ESTADOS
% ==============================
\section{Diagrama de estados}
El diagrama representa una Máquina de Estados Finitos (FSM) síncrona de tipo Moore diseñada para el control de acceso mediante una clave binaria de 5 bits (\texttt{1-0-1-1-0}). Cada transición entre estados se valida ante un flanco activo de la señal \texttt{ENTER} evaluando la entrada \texttt{X}. La arquitectura implementa un mecanismo de seguridad basado en dos ramas paralelas: la rama de aciertos (\texttt{OK1} a \texttt{OK4}) y la rama de error arrastrado (\texttt{E1} a \texttt{E4}). Si el usuario introduce un bit incorrecto en cualquier etapa, el flujo se desvía a la rama de error, la cual continúa avanzando en cada pulsación de \texttt{ENTER} indiferente al valor de \texttt{X}, reflejando exactamente el mismo patrón de progreso en los LEDs (\texttt{0001} a \texttt{0100}) para no alertar al usuario sobre en qué posición ocurrió el fallo. Al ingresar el quinto bit, si la secuencia fue correcta se alcanza el estado \texttt{ESTADO\_OK} activando \texttt{ACCESO\_OK}; en caso contrario, se evalúa el contador de fallos (\texttt{cont\_intentos}): si es menor a 2 se transita a \texttt{NOK} permitiendo reiniciar el proceso con otro pulso de \texttt{ENTER}, pero si se alcanza el tercer fallo acumulado ($\texttt{cont\_intentos} = 2$), el sistema pasa de manera irreversible al estado \texttt{BLOQUEO} ($\texttt{BLOQUEADO} = \text{'1'}$), del cual solo es posible salir mediante una señal global de \texttt{RESET}.

\begin{figure}[H]
    \centering
    \includegraphics[width=1\linewidth]{diagram.png}
    \caption{Diagrama de estados de la FSM incorporando rama de error.}
    \label{fig:diagrama_estados}
\end{figure}

% ==============================
% 8. TABLA DE TRANSICIÓN
% ==============================
\section{Tabla de transición}

\begin{table}[H]
\centering
\caption{Tabla de transiciones y salidas de la FSM.}
\label{tab:fsm_transiciones}
\small
\resizebox{\textwidth}{!}{%
\begin{tabular}{|l|c|l|c|c|c|c|}
\hline
\textbf{Estado} & \textbf{Entrada} & \textbf{Estado} & \multicolumn{4}{c|}{\textbf{Salidas de Moore}} \\ \cline{4-7}
\textbf{Actual} & \textbf{$X$} & \textbf{Siguiente} & \texttt{ESTADO\_LED} & \texttt{ACCESO\_OK} & \texttt{ERROR} & \texttt{BLOQUEADO} \\ \hline
\hline
\texttt{INICIO}    & 1 & \texttt{OK1}       & \texttt{0000} & 0 & 0 & 0 \\
                   & 0 & \texttt{E1}        & \texttt{0000} & 0 & 0 & 0 \\ \hline
\texttt{OK1}       & 0 & \texttt{OK2}       & \texttt{0001} & 0 & 0 & 0 \\
                   & 1 & \texttt{E2}        & \texttt{0001} & 0 & 0 & 0 \\ \hline
\texttt{OK2}       & 1 & \texttt{OK3}       & \texttt{0010} & 0 & 0 & 0 \\
                   & 0 & \texttt{E3}        & \texttt{0010} & 0 & 0 & 0 \\ \hline
\texttt{OK3}       & 1 & \texttt{OK4}       & \texttt{0011} & 0 & 0 & 0 \\
                   & 0 & \texttt{E4}        & \texttt{0011} & 0 & 0 & 0 \\ \hline
\texttt{OK4}       & 0 & \texttt{ESTADO\_OK}& \texttt{0100} & 0 & 0 & 0 \\
                   & 1 & \texttt{NOK} ($cont < 2$) / \texttt{BLOQUEO} ($cont = 2$) & \texttt{0100} & 0 & 0 & 0 \\ \hline
\texttt{E1}        & X & \texttt{E2}        & \texttt{0001} & 0 & 0 & 0 \\ \hline
\texttt{E2}        & X & \texttt{E3}        & \texttt{0010} & 0 & 0 & 0 \\ \hline
\texttt{E3}        & X & \texttt{E4}        & \texttt{0011} & 0 & 0 & 0 \\ \hline
\texttt{E4}        & X & \texttt{NOK} ($cont < 2$) / \texttt{BLOQUEO} ($cont = 2$) & \texttt{0100} & 0 & 0 & 0 \\ \hline
\texttt{ESTADO\_OK}& - & \texttt{ESTADO\_OK}& \texttt{0101} & \textbf{1} & 0 & 0 \\ \hline
\texttt{NOK}       & - & \texttt{INICIO}    & \texttt{1110} & 0 & \textbf{1} & 0 \\ \hline
\texttt{BLOQUEO}   & - & \texttt{BLOQUEO}   & \texttt{1111} & 0 & 0 & \textbf{1} \\ \hline
\end{tabular}%
}
\end{table}

% ==============================
% 9. CÓDIGO VHDL
% ==============================
\section{Código VHDL}
El código fuente consta de un sincronizador de pulsos (\texttt{detector\_flanco.vhd}), la lógica de estados (\texttt{fsm\_acceso.vhd}) y el módulo de integración top-level (\texttt{lab03\_top.vhd}). Los listados completos del código se encuentran adjuntos en el apartado de Anexos.

% ==============================
% 10. TESTBENCH
% ==============================
\section{Testbench}
El archivo \texttt{tb\_fsm\_acceso.vhd} contiene el banco de pruebas. Utiliza un reloj simulado a 100 MHz y evalúa seis casos de prueba obligatorios mediante aserciones (\texttt{assert}), abarcando ingresos exitosos, fallos múltiples y recuperación por bloqueo tras tres intentos. El código del banco de pruebas está detallado en los Anexos.

% ==============================
% 11. SIMULACIÓN EN VIVADO
% ==============================
\section{Simulación en Vivado}
La siguiente captura evidencia el comportamiento de las señales internas (\texttt{CLK}, \texttt{RESET}, \texttt{X}, \texttt{ENTER}, \texttt{pulso\_enter}, \texttt{estado\_actual}, \texttt{estado\_siguiente}, \texttt{INTENTOS}, \texttt{ACCESO\_OK}, \texttt{ERROR}, \texttt{BLOQUEADO}) en el entorno de simulación de Vivado, validando la transición de estados sin la presencia de temporizaciones no deseadas.

\begin{figure}[H]
    \centering
    \includegraphics[width=1\linewidth]{simulacion.png}
    \caption{Simulación en Vivado.}
    \label{fig:simulacion}
\end{figure}


% ==============================
% 12. ARCHIVO DE RESTRICCIONES
% ==============================
\section{Archivo de restricciones}
El archivo \texttt{lab03.xdc} relaciona los puertos lógicos declarados en la entidad top-level con los pines físicos de la FPGA Nexys 4 DDR, asegurando que la botonera, los interruptores y los LEDs respondan adecuadamente a las señales de la tarjeta. El código se encuentra en la sección de Anexos.

\section{Resultados obtenidos}
La implementación física en la placa FPGA Nexys 4 DDR permitió comprobar el correcto funcionamiento de la máquina de estados. El sistema validó con éxito la clave \texttt{10110}, registró adecuadamente los intentos fallidos y bloqueó el sistema al alcanzar el tercer error.

\begin{figure}[H]
    \centering
    \includegraphics[width=0.8\linewidth]{estado_ok.jpg}
    \caption{Implementación física del sistema de control evidenciando la validación exitosa de la clave (ESTADO\_OK). Se observa el LED1 (ACCESO\_OK) encendido y el código de estado 0101.}
    \label{fig:estado_ok}
\end{figure}

\begin{figure}[H]
    \centering
    \includegraphics[width=0.8\linewidth]{nok_1er_intento.jpg}
    \caption{Implementación física evidenciando el estado de error (NOK) tras registrar el primer intento fallido. Se observa el LED2 (ERROR) encendido, el contador en 1 (LED4) y el código de estado 1110.}
    \label{fig:nok_1er_intento}
\end{figure}

\begin{figure}[H]
    \centering
    \includegraphics[width=0.8\linewidth]{nok_2do_intento.jpg}
    \caption{Implementación física evidenciando el estado de error (NOK) tras registrar el segundo intento fallido. Se observa el LED2 (ERROR) encendido, el contador en 2 (LED5) y el código de estado 1110.}
    \label{fig:nok_2do_intento}
\end{figure}

\begin{figure}[H]
    \centering
    \includegraphics[width=0.8\linewidth]{estado_bloqueo.jpg}
    \caption{Implementación física evidenciando el estado de BLOQUEO permanente tras tres errores. Se observa la señal de bloqueo activada (LED3), el contador de intentos al máximo (LED4 y LED5) y el código visual de estado bloqueado (LEDs 6 al 9 encendidos).}
    \label{fig:estado_bloqueo}
\end{figure}
% ==============================
% 13. ANÁLISIS DE ERRORES
% ==============================
\section{Análisis de errores}
Durante el desarrollo lógico y la simulación del sistema, se enfrentó y corrigió un desafío principal:
\begin{itemize}
    \item \textbf{Seguridad de la FSM inicial:} En el primer modelo lógico, si se ingresaba un bit incorrecto, la FSM saltaba instantáneamente al estado de error. Esto revelaba de manera prematura en qué parte de la clave se había equivocado el usuario. Se resolvió rediseñando el diagrama para incluir estados de error (\texttt{E1} a \texttt{E4}) que permiten completar la inserción de los 5 bits simulando una secuencia normal antes de encender la señal de \texttt{ERROR}.
\end{itemize}

% ==============================
% 14. RESPUESTAS A PREGUNTAS
% ==============================
\section{Respuestas a preguntas}
De acuerdo al análisis requerido en el laboratorio, se presentan las siguientes respuestas:
\begin{enumerate}
    \item \textbf{¿Cuál es la diferencia entre un circuito combinacional y uno secuencial?}\\
    Un circuito combinacional genera salidas que dependen exclusivamente de sus entradas actuales. Un circuito secuencial, por el contrario, posee elementos de memoria interna (como biestables), por lo que sus salidas dependen tanto de las entradas actuales como del estado anterior del sistema.
    
    \item \textbf{¿Qué función cumple el registro de estado?}\\
    Es el componente de hardware sincronizado por reloj encargado de almacenar y retener el estado actual en el que se encuentra la máquina hasta que ocurra la próxima transición.
    
    \item \textbf{¿Cuál es la diferencia entre una máquina de Moore y una de Mealy?}\\
    En la máquina de Moore, los valores de las salidas están determinados únicamente por el estado actual. En la máquina de Mealy, las salidas dependen tanto del estado actual como del valor de las entradas en ese instante.
    
    \item \textbf{¿Por qué se utilizan tipos enumerados para representar estados?}\\
    Permiten que el código sea legible y semántico para el programador (por ejemplo, usar \texttt{INICIO} en lugar de \texttt{"000"}). Además, confieren libertad al sintetizador de Vivado para elegir la codificación binaria más óptima en hardware (como \textit{One-Hot}).
    
    \item \textbf{¿Por qué se deben asignar valores por defecto?}\\
    Para prevenir la inferencia de \textit{latches} (memorias no deseadas) durante el proceso de síntesis. Si una señal no recibe valor en alguna de las ramas de condición lógica, el sistema mantendrá su valor anterior generando hardware asíncrono indeseado.
    
    \item \textbf{¿Por qué no se recomienda usar un pulsador directamente como reloj?}\\
    Porque la señal mecánica de un pulsador contiene ruido que el sistema interpreta como múltiples flancos rápidos y provoca problemas severos de \textit{clock skew}. Siempre se debe usar el reloj de cuarzo de alta frecuencia de la placa.
    
    \item \textbf{¿Qué efecto produce el rebote de un pulsador?}\\
    Produce fluctuaciones de voltaje transitorias. A frecuencias de megahercios, la FPGA interpretará un solo toque humano como si el botón hubiera sido presionado docenas de veces, alterando por completo la secuencia de estados.
    
    \item \textbf{¿Qué función cumple el reset?}\\
    Devuelve la máquina de estados a un punto de partida inicial seguro, limpiando los registros y borrando contadores para asegurar un funcionamiento predecible desde el momento del encendido.
    
    \item \textbf{¿Cómo se implementan los estados dentro de la FPGA?}\\
    Los estados se almacenan mediante Flip-Flops (elementos de memoria) dentro de los \textit{Slices} de la FPGA, mientras que la lógica de transición y decisión se mapea sobre las Tablas de Búsqueda (LUTs).
    
    \item \textbf{¿Qué relación existe entre una FSM y la unidad de control de un procesador?}\\
    La unidad de control de un procesador es en su esencia una FSM compleja. En lugar de procesar una clave, procesa instrucciones máquina (como Fetch, Decode, Execute) dictando la secuencia de apertura de buses, registros y operaciones de la ALU ciclo a ciclo.
\end{enumerate}

% ==============================
% 15. CONCLUSIONES
% ==============================
\section{Conclusiones}
La implementación de la máquina de estados finitos demostró la eficacia de VHDL para abstraer procesos lógicos y llevarlos a componentes físicos en la FPGA Nexys 4 DDR. Durante el laboratorio se evidenció de manera práctica la importancia de mantener un diseño puramente síncrono manejado por el reloj principal, utilizando el pulsador físico solo como una entrada de datos condicionada mediante un bloque detector de flancos. Adicionalmente, se logró integrar lógica de seguridad a nivel de hardware, añadiendo estados ocultos para evitar revelar la posición del fallo durante el ingreso de la clave, cumpliendo íntegramente con los requisitos de simulación, conteo de intentos y bloqueo permanente.

\newpage
% ==============================
% REFERENCIAS
% ==============================
\addcontentsline{toc}{section}{Referencias}
\begin{thebibliography}{9}

\bibitem{nexys4ddr}
Digilent Inc., \emph{Nexys 4 DDR FPGA Board Reference Manual}, DOC\#: 502-292. [En línea]. Disponible en: \url{https://digilent.com/reference/programmable-logic/nexys-4-ddr/reference-manual}

\bibitem{artix7}
Xilinx, Inc., \emph{7 Series FPGAs Data Sheet: Overview}, DS180.

\bibitem{ieee1076}
IEEE, \emph{IEEE Standard VHDL Language Reference Manual}, IEEE Std 1076-2008, 2009.

\bibitem{ieee1164}
IEEE, \emph{IEEE Standard Multivalue Logic System for VHDL Model Interoperability (Std\_logic\_1164)}, IEEE Std 1164-1993, 1993.

\bibitem{ieee1076_3}
IEEE, \emph{IEEE Standard VHDL Synthesis Packages}, IEEE Std 1076.3-1997, 1997.

\bibitem{harris}
D. M. Harris y S. L. Harris, \emph{Digital Design and Computer Architecture}, 2.ª ed. Morgan Kaufmann, 2012.

\bibitem{pedroni}
V. A. Pedroni, \emph{Circuit Design and Simulation with VHDL}, 2.ª ed. MIT Press, 2010.

\end{thebibliography}

\newpage
% ==============================
% 7. ANEXOS
% ==============================
\section{Anexos}

\subsection{Repositorio del proyecto}

\begin{itemize}
    \item \url{https://github.com/JoakinMacAuliffe/Lab3_Arqui}
\end{itemize}

\subsection{Código fuente VHDL}

\subsubsection{\texttt{detector\_flanco.vhd} (Módulo sincronizador y detector de flancos)}
\begin{lstlisting}[language=VHDL, caption={detector\_flanco.vhd}]
library IEEE;
use IEEE.STD_LOGIC_1164.ALL;

entity detector_flanco is
    Port (
        CLK : in STD_LOGIC;
        BOTON : in STD_LOGIC;
        PULSO : out STD_LOGIC
    );
end detector_flanco;

architecture Behavioral of detector_flanco is
    signal q0, q1 : std_logic := '0';
begin
    process(CLK)
    begin
        if rising_edge(CLK) then
            q0 <= BOTON;
            q1 <= q0;
        end if;
    end process;
    
    PULSO <= q0 and not q1;
end Behavioral;
\end{lstlisting}

\subsubsection{\texttt{fsm\_acceso.vhd} (La máquina de estados modificada)}
\begin{lstlisting}[language=VHDL, caption={fsm\_acceso.vhd}]
library IEEE;
use IEEE.STD_LOGIC_1164.ALL;
use IEEE.NUMERIC_STD.ALL;

entity fsm_acceso is
    Port ( 
        CLK       : in STD_LOGIC;
        RESET     : in STD_LOGIC;
        X         : in STD_LOGIC;
        ENTER     : in STD_LOGIC;
        ACCESO_OK : out STD_LOGIC;
        ERROR     : out STD_LOGIC;
        BLOQUEADO : out STD_LOGIC;
        INTENTOS  : out STD_LOGIC_VECTOR(1 downto 0);
        ESTADO_LED: out STD_LOGIC_VECTOR(3 downto 0)
    );
end fsm_acceso;

architecture Behavioral of fsm_acceso is
    -- Tipo enumerado de estados: rama de aciertos (OK1-OK4) y rama de error (E1-E4)
    type tipo_estado is (INICIO, OK1, OK2, OK3, OK4, ESTADO_OK, E1, E2, E3, E4, NOK, BLOQUEO);
    signal estado_actual, estado_siguiente : tipo_estado;
    signal cont_intentos : unsigned(1 downto 0) := "00";
begin

    -- Registro de estado y contador de intentos (Sincrono)
    process (CLK)
    begin
        if rising_edge(CLK) then
            if RESET = '1' then
                estado_actual <= INICIO;
                cont_intentos <= "00";
            elsif ENTER = '1' then
                if estado_actual = BLOQUEO then
                    estado_actual <= BLOQUEO;
                else
                    estado_actual <= estado_siguiente;
                    
                    -- El contador solo se incrementa al ingresar el 5to bit (transicion a NOK o BLOQUEO)
                    if (estado_siguiente = NOK or estado_siguiente = BLOQUEO) and 
                       (estado_actual = E4 or estado_actual = OK4) then
                        if cont_intentos < 3 then
                            cont_intentos <= cont_intentos + 1;
                        end if;
                    end if;
                end if;
            end if;
        end if;
    end process;

    -- Logica de transicion de estados (combinacional)
    process (estado_actual, X, cont_intentos)
    begin
        estado_siguiente <= estado_actual;
        case estado_actual is
            -- Rama de aciertos
            when INICIO => if X = '1' then estado_siguiente <= OK1; else estado_siguiente <= E1; end if;
            when OK1    => if X = '0' then estado_siguiente <= OK2; else estado_siguiente <= E2; end if;
            when OK2    => if X = '1' then estado_siguiente <= OK3; else estado_siguiente <= E3; end if;
            when OK3    => if X = '1' then estado_siguiente <= OK4; else estado_siguiente <= E4; end if;
            when OK4    => 
                if X = '0' then 
                    estado_siguiente <= ESTADO_OK; 
                else 
                    if cont_intentos = 2 then estado_siguiente <= BLOQUEO; else estado_siguiente <= NOK; end if;
                end if;

            -- Rama de error arrastrado
            when E1 => estado_siguiente <= E2; -- Ignora X, arrastra el error
            when E2 => estado_siguiente <= E3;
            when E3 => estado_siguiente <= E4;
            when E4 => 
                if cont_intentos = 2 then estado_siguiente <= BLOQUEO; else estado_siguiente <= NOK; end if;

            -- Estados finales (NOK vuelve a INICIO con ENTER; ESTADO_OK y BLOQUEO se mantienen)
            when ESTADO_OK => estado_siguiente <= ESTADO_OK;
            when NOK       => estado_siguiente <= INICIO; -- Al presionar ENTER despues del error, vuelve al inicio
            when BLOQUEO   => estado_siguiente <= BLOQUEO;
            when others    => estado_siguiente <= INICIO;
        end case;
    end process;

    -- Logica de salida (Moore): depende solo del estado actual
    process (estado_actual, cont_intentos)
    begin
        ACCESO_OK <= '0';
        ERROR <= '0';
        BLOQUEADO <= '0';
        ESTADO_LED <= "0000";

        case estado_actual is
            when INICIO => ESTADO_LED <= "0000";
            
            -- Los estados E muestran el mismo codigo de LEDs que los OK, para no revelar la posicion del error
            when OK1 | E1 => ESTADO_LED <= "0001";
            when OK2 | E2 => ESTADO_LED <= "0010";
            when OK3 | E3 => ESTADO_LED <= "0011";
            when OK4 | E4 => ESTADO_LED <= "0100";
            
            when ESTADO_OK => 
                ACCESO_OK <= '1';
                ESTADO_LED <= "0101";
                
            -- El error solo se revela al final (NOK)
            when NOK => 
                ERROR <= '1';
                ESTADO_LED <= "1110";
                
            when BLOQUEO => 
                BLOQUEADO <= '1';
                ESTADO_LED <= "1111";
            when others => null;
        end case;
        
        INTENTOS <= std_logic_vector(cont_intentos);
    end process;

end Behavioral;
\end{lstlisting}

\subsubsection{\texttt{lab03\_top.vhd} (El módulo principal que integra el diseño)}
\begin{lstlisting}[language=VHDL, caption={lab03\_top.vhd}]
library IEEE;
use IEEE.STD_LOGIC_1164.ALL;

entity lab03_top is
    Port (
        CLK       : in STD_LOGIC;
        RESET     : in STD_LOGIC;
        X         : in STD_LOGIC;
        ENTER     : in STD_LOGIC;
        X_OUT     : out STD_LOGIC;
        ACCESO_OK : out STD_LOGIC;
        ERROR     : out STD_LOGIC;
        BLOQUEADO : out STD_LOGIC;
        INTENTOS  : out STD_LOGIC_VECTOR(1 downto 0);
        ESTADO_LED: out STD_LOGIC_VECTOR(3 downto 0)
    );
end lab03_top;

architecture Behavioral of lab03_top is
    signal pulso_enter : STD_LOGIC;
begin
    X_OUT <= X;

    Inst_detector: entity work.detector_flanco
        port map(
            CLK => CLK,
            BOTON => ENTER,
            PULSO => pulso_enter
        );

    Inst_fsm: entity work.fsm_acceso
        port map(
            CLK => CLK,
            RESET => RESET,
            X => X,
            ENTER => pulso_enter,
            ACCESO_OK => ACCESO_OK,
            ERROR => ERROR,
            BLOQUEADO => BLOQUEADO,
            INTENTOS => INTENTOS,
            ESTADO_LED => ESTADO_LED
        );
end Behavioral;
\end{lstlisting}

\subsection{Testbench (Banco de pruebas)}

\subsubsection{\texttt{tb\_fsm\_acceso.vhd} (El código de simulación con los casos de prueba corregidos)}
\begin{lstlisting}[language=VHDL, caption={tb\_fsm\_acceso.vhd}]
library IEEE;
use IEEE.STD_LOGIC_1164.ALL;

entity tb_fsm_acceso is
end tb_fsm_acceso;

architecture behavior of tb_fsm_acceso is
    signal CLK : std_logic := '0';
    signal RESET : std_logic := '0';
    signal X : std_logic := '0';
    signal ENTER : std_logic := '0';
    
    signal X_OUT : std_logic;
    signal ACCESO_OK : std_logic;
    signal ERROR : std_logic;
    signal BLOQUEADO : std_logic;
    signal INTENTOS : std_logic_vector(1 downto 0);
    signal ESTADO_LED : std_logic_vector(3 downto 0);
begin
    uut: entity work.lab03_top PORT MAP (
          CLK => CLK, RESET => RESET, X => X, ENTER => ENTER,
          X_OUT => X_OUT, ACCESO_OK => ACCESO_OK, ERROR => ERROR,
          BLOQUEADO => BLOQUEADO, INTENTOS => INTENTOS, ESTADO_LED => ESTADO_LED
        );

    CLK <= not CLK after 5 ns;

    stim_proc: process
        procedure pulso_btn is
        begin
            ENTER <= '1'; wait for 20 ns;
            ENTER <= '0'; wait for 20 ns;
        end procedure;
    begin
        RESET <= '1'; wait for 20 ns;
        RESET <= '0'; wait for 20 ns;
        
        -- CASO 1: Secuencia Correcta 10110
        X<='1'; pulso_btn; X<='0'; pulso_btn; X<='1'; pulso_btn; X<='1'; pulso_btn; X<='0'; pulso_btn;
        wait for 10 ns;
        assert (ACCESO_OK = '1') report "Fallo Caso 1: Acceso no autorizado" severity FAILURE;
        
        -- CASO 2: Error, Intento 1 (00000)
        RESET <= '1'; wait for 20 ns; RESET <= '0'; wait for 20 ns;
        X<='0'; pulso_btn; X<='0'; pulso_btn; X<='0'; pulso_btn; X<='0'; pulso_btn; X<='0'; pulso_btn;
        wait for 10 ns;
        assert (INTENTOS = "01") report "Fallo Caso 2: Intento 1 no registrado" severity FAILURE;
        
        -- Pulso extra para salir del estado NOK y volver a INICIO sin borrar el contador
        pulso_btn;
        
        -- CASO 3: Error, Intento 2 (11111)
        X<='1'; pulso_btn; X<='1'; pulso_btn; X<='1'; pulso_btn; X<='1'; pulso_btn; X<='1'; pulso_btn;
        wait for 10 ns;
        assert (INTENTOS = "10") report "Fallo Caso 3: Intento 2 no registrado" severity FAILURE;

        -- Pulso extra para salir del estado NOK y volver a INICIO sin borrar el contador
        pulso_btn;

        -- CASO 4: Error, Intento 3 y BLOQUEO (10111)
        X<='1'; pulso_btn; X<='0'; pulso_btn; X<='1'; pulso_btn; X<='1'; pulso_btn; X<='1'; pulso_btn;
        wait for 10 ns;
        assert (BLOQUEADO = '1') report "Fallo Caso 4: Bloqueo no activado" severity FAILURE;

        -- CASO 5: Despues del bloqueo (10110) - NO debe autorizar
        X<='1'; pulso_btn; X<='0'; pulso_btn; X<='1'; pulso_btn; X<='1'; pulso_btn; X<='0'; pulso_btn;
        wait for 10 ns;
        assert (ACCESO_OK = '0' and BLOQUEADO = '1') report "Fallo Caso 5: Se autorizo estando bloqueado" severity FAILURE;

        -- CASO 6: RESET y clave correcta (10110)
        RESET <= '1'; wait for 20 ns; RESET <= '0'; wait for 20 ns;
        X<='1'; pulso_btn; X<='0'; pulso_btn; X<='1'; pulso_btn; X<='1'; pulso_btn; X<='0'; pulso_btn;
        wait for 10 ns;
        assert (ACCESO_OK = '1') report "Fallo Caso 6: No autorizo tras el reset" severity FAILURE;

        wait;
    end process;
end behavior;
\end{lstlisting}

\subsection{Archivo XDC (Restricciones)}

\subsubsection{\texttt{lab03.xdc} (La asignación física de pines para la FPGA)}
\begin{lstlisting}[language=tcl, caption={lab03.xdc}]
# Reloj
set_property -dict { PACKAGE_PIN E3    IOSTANDARD LVCMOS33 } [get_ports { CLK }];
create_clock -add -name sys_clk_pin -period 10.00 -waveform {0 5} [get_ports {CLK}];

# Interruptor SW0 (Entrada X)
set_property -dict { PACKAGE_PIN J15   IOSTANDARD LVCMOS33 } [get_ports { X }];

# Botones (ENTER y RESET)
set_property -dict { PACKAGE_PIN N17   IOSTANDARD LVCMOS33 } [get_ports { ENTER }]; # BTNC
set_property -dict { PACKAGE_PIN M18   IOSTANDARD LVCMOS33 } [get_ports { RESET }]; # BTNU

# LEDs de Salida
set_property -dict { PACKAGE_PIN H17   IOSTANDARD LVCMOS33 } [get_ports { X_OUT }]; # LED0
set_property -dict { PACKAGE_PIN K15   IOSTANDARD LVCMOS33 } [get_ports { ACCESO_OK }]; # LED1
set_property -dict { PACKAGE_PIN J13   IOSTANDARD LVCMOS33 } [get_ports { ERROR }]; # LED2
set_property -dict { PACKAGE_PIN N14   IOSTANDARD LVCMOS33 } [get_ports { BLOQUEADO }]; # LED3

# LEDs Intentos (LED4 y LED5)
set_property -dict { PACKAGE_PIN R18   IOSTANDARD LVCMOS33 } [get_ports { INTENTOS[0] }];
set_property -dict { PACKAGE_PIN V17   IOSTANDARD LVCMOS33 } [get_ports { INTENTOS[1] }];

# LEDs Estado (LED6 a LED9)
set_property -dict { PACKAGE_PIN U17   IOSTANDARD LVCMOS33 } [get_ports { ESTADO_LED[0] }];
set_property -dict { PACKAGE_PIN U16   IOSTANDARD LVCMOS33 } [get_ports { ESTADO_LED[1] }];
set_property -dict { PACKAGE_PIN V16   IOSTANDARD LVCMOS33 } [get_ports { ESTADO_LED[2] }];
set_property -dict { PACKAGE_PIN T15   IOSTANDARD LVCMOS33 } [get_ports { ESTADO_LED[3] }];
\end{lstlisting}

\end{document}